%%%%%%%%%%%%%%%%%%%%%%%%%%%% Document Setup %%%%%%%%%%%%%%%%%%
\documentclass[10pt]{article}
\usepackage[T2A,T1]{fontenc}
\usepackage[utf8]{inputenc}
\usepackage[russian,english]{babel}
\usepackage{CJKutf8}
\AtBeginDvi{\input{zhwinfonts}}
%\usepackage{polyglossia}
\makeatletter
\newlength{\bibhang}
\setlength{\bibhang}{1em}
\newlength{\bibsep}
 {\@listi \global\bibsep\itemsep \global\advance\bibsep by\parsep}
\makeatother
\reversemarginpar
\usepackage{geometry}\geometry{a4paper,total={170mm,257mm},left=1.18in,right=1.18in,top=0.79in,bottom=0.79in}

\setlength{\parindent}{0in}
\usepackage{enumerate} 
\usepackage{etaremune}
\usepackage{calc}
\usepackage{bbding}
\usepackage{paralist}
\usepackage{fancyhdr,lastpage}
\pagestyle{fancy}
%\pagestyle{empty} % Uncomment this to get rid of page numbers
\fancyhf{}\renewcommand{\headrulewidth}{0pt}
\fancyfootoffset{\marginparsep+\marginparwidth}
\newlength{\footpageshift}
\setlength{\footpageshift}
          {0.5\textwidth+1.\marginparsep+0.5\marginparwidth-2in}
\lfoot{\hspace{\footpageshift}%
       \parbox{4in}{\, \hfill %
                    \arabic{page} of \protect\pageref*{LastPage} % +LP
% \arabic{page} % -LP
                    \hfill \,}}
\usepackage{color,hyperref}
\usepackage{graphics}
\usepackage{wrapfig}
\usepackage{longtable}

\usepackage[usenames,dvipsnames]{xcolor}
\definecolor{darkblue}{rgb}{0.0,0.0,0.3}
\definecolor{mygreen}{RGB}{44,204,148}
\definecolor{darkgreen}{RGB}{13,107,27}
\definecolor{lightgreen}{RGB}{0,170,37}
\definecolor{provinces}{RGB}{11,100,132}
\definecolor{mypurple}{RGB}{77,24,137}
\definecolor{haiyan}{RGB}{8,150,178}
\definecolor{anticred}{RGB}{165,62,11}
\definecolor{amber}{rgb}{1.0, 0.75, 0.0}
\definecolor{americanrose}{rgb}{1.0, 0.01, 0.24}
\definecolor{applegreen}{rgb}{0.55, 0.71, 0.0}
\definecolor{bananayellow}{rgb}{1.0, 0.88, 0.21}
\definecolor{beaublue}{rgb}{0.74, 0.83, 0.9}
\definecolor{bittersweet}{rgb}{1.0, 0.44, 0.37}
\definecolor{blue-violet}{rgb}{0.54, 0.17, 0.89}
\definecolor{brilliantlavender}{rgb}{0.96, 0.73, 1.0}
\definecolor{brilliantrose}{rgb}{1.0, 0.33, 0.64}
\definecolor{brightturquoise}{rgb}{0.03, 0.91, 0.87}

\hypersetup{colorlinks,breaklinks,
            linkcolor=darkblue,urlcolor=darkblue,
            anchorcolor=darkblue,citecolor=darkblue}
%%%%%%%%%%%%%%%%%%%%%%%% End Document Setup %%%%%%%%%%%%%%%%%%%
%%%%%%%%%%%%%%%%%%%%%%%%%%% Helper Commands %%%%%%%%%

\newcommand{\makeheading}[1]%
        {\hspace*{-\marginparsep minus \marginparwidth}%
         \begin{minipage}[t]{\textwidth+\marginparwidth+\marginparsep}%
                {\large \bfseries #1}\\[-0.15\baselineskip]%
                 \rule{\columnwidth}{1pt}%
         \end{minipage}}
\renewcommand{\section}[2]%
        {\pagebreak[2]\vspace{1.3\baselineskip}%
         \phantomsection\addcontentsline{toc}{section}{#1}%
         \hspace{0in}%
         \marginpar{
         \raggedright \scshape #1}#2}
\newenvironment{outerlist}[1][\enskip\textbullet]%
        {\begin{itemize}[#1]}{\end{itemize}%
         \vspace{-.6\baselineskip}}
\newenvironment{lonelist}[1][\enskip\textbullet]%
        {\vspace{-\baselineskip}\begin{list}{#1}{%
        \setlength{\partopsep}{0pt}%
        \setlength{\topsep}{0pt}}}
        {\end{list}\vspace{-.6\baselineskip}}
\newenvironment{innerlist}[1][\enskip\textbullet]%
        {\begin{compactitem}[#1]}{\end{compactitem}}

\newenvironment{loneinnerlist}[1][\enskip\textbullet]%
        {\vspace{-\baselineskip}\begin{compactitem}[#1]}
        {\end{compactitem}\vspace{-.6\baselineskip}}
\newcommand{\blankline}{\quad\pagebreak[2]}
\newcommand\doilink[1]{\href{http://dx.doi.org/#1}{#1}}
\newcommand\doi[1]{doi:\doilink{#1}}
\providecommand*\url[1]{\href{#1}{#1}}
\renewcommand*\url[1]{\href{#1}{\texttt{#1}}}

\usepackage{tikz}
\usetikzlibrary{chains}
\usepackage{pgfplots}
\usepackage{pgfmath}
\usetikzlibrary{matrix}
\usetikzlibrary{arrows,decorations.pathmorphing,backgrounds,positioning,fit,petri}
\usepgfmodule{decorations}
\usepgflibrary{decorations.shapes} 
\usetikzlibrary[decorations.shapes] 
\usetikzlibrary{mindmap}
\usepgflibrary{patterns}
\usetikzlibrary[patterns]
\usetikzlibrary{petri}
\usepgflibrary{plothandlers}
\usetikzlibrary{plothandlers}
\usepackage{pgfpages}
\usepgflibrary{arrows}
\usetikzlibrary{arrows}
\usetikzlibrary{trees}
\usetikzlibrary{topaths}
\usetikzlibrary{positioning}
\usetikzlibrary{fit}
\usepgflibrary{shapes.misc}
\usetikzlibrary{shapes.misc}
\usetikzlibrary{spy}
\usepgflibrary{shadings}
\usetikzlibrary{shadings}
\usepgfmodule{shapes}
\usepgfplotslibrary{external}
\usetikzlibrary{pgfplots.external}
\usepgfplotslibrary{colorbrewer}
\usepgfplotslibrary{patchplots}
\usetikzlibrary{plotmarks}
\usepackage{color} 
\usepackage{xcolor} 
%\pgfpagesuselayout{resize to}[a4paper,border shrink=20mm]
\pgfplotsset{colormap={cool}{rgb255(0cm)=(255,255,255); rgb255(1cm)=(0,128,255); rgb255(2cm)=(255,0,255)}}
  \pgfplotsset{colormap={violet}{rgb255=(25,25,122) color=(white) rgb255=(238,140,238)}}
  \pgfplotsset{
    colormap={bluered}{rgb255(0cm)=(0,0,180); rgb255(1cm)=(0,255,255); rgb255(2cm)=(100,255,0);
rgb255(3cm)=(255,255,0); rgb255(4cm)=(255,0,0); rgb255(5cm)=(128,0,0)}}
 \pgfplotsset{/pgfplots/colormap={winter}{rgb255=(0,0,255) rgb255=(0,255,128)}}
 \pgfplotsset{
    /pgfplots/colormap={spring}{rgb255=(255,0,255) rgb255=(255,255,0)}}
\definecolor{bisque}{rgb}{1.0, 0.89, 0.77}

\usepgfplotslibrary{ternary} % <-- use 'ternary' library for plotting triangle diagrams
\usetikzlibrary{pgfplots.ternary} % <-- we explicitly prescribe that we will use the 'ternary' libraries in the picture environment  
 
\begin{document} % <-- start the document
\selectlanguage{russian}
\begin{tikzpicture} [scale=1]% <--  start subplot 1
\begin{ternaryaxis}% <-- start ternary diagram
[title={\textbf{Granulometric composition of the soil}. Fractions: dust, clay, sand}, % <-- set up subplot parameters in square brackets
 legend cell align=left,% <-- align legend column on the left
 legend pos=outer north east,
 grid=both, % <--  draw a detailed two-sided grid
 legend entries={sand, loamy sand, sandy loam, sandy clay loam, sandy clay, clay, clay loam, silty clay, silty clay loam, clay loam, silty clay loam, fine loam, silty clay}, 
    xlabel=\footnotesize{Clay (<0.002 mm)}, % <-- x axis - Name of fraction: dust. Specify a small font size: \ footnotesize
    ylabel=\footnotesize{Dust (0.002-0.05 mm)}, % <-- y axis - Name of fraction: clay
    zlabel=\footnotesize{Sand (0.05-2 mm)},% <-- z axis Name of fraction: sand
    label style={sloped},% <-- indicate tilt annotations along the axises
    minor tick num=2] % <-- indicate how many minor ticks are between the major ones on the coordinate grid (in this case, 2).
    
   \addplot3 coordinates {  % <-- start series of objects, enter coordinates in ternary mode
         (0.60,    0.40,    0.00)
        (0.40,    0.60,    0.00)
        (0.40,    0.60,    0.20)
};  \addlegendentry{sandy clay}% <-- add object into legend   
      \addplot3 coordinates {
        (0.20,    0.80,    0.80)
        (0.20,    0.80,    0.60)
        (0.10,    0.90,    0.50)
        (0.10,    0.90,    0.40)
       (0.00,    1.00,    0.50)  
        (0.00,    1.00,    0.60)       
        (0.10,    0.90,    0.90)
 };  \addlegendentry{sandy loam}

   \addplot3 coordinates { 
         (0.0,    1.00,    0.60) 
        (0.0,    1.00,    0.80)
        (0.10,    0.95,    0.95)
         (0.15,    0.85,    0.85)
};  \addlegendentry{loamy sand}

   \addplot3 coordinates { 
      (0.0,    1.00,    1.00) 
      (0.10,    0.85,    0.95) 
      (0.00,    1.00,    0.80)
};  \addlegendentry{sand}

   \addplot3 coordinates { 
        (1.00,    0.00,    0.00) 
        (0.60,    0.40,    0.00)        
        (0.40,    0.60,    0.20)    
        (0.40,    0.60,    0.40)
        (0.60,    0.40,    0.40)     
}; \addlegendentry{clay} 
  
     \addplot3 coordinates {  
        (0.40,    0.60,    0.40)
        (0.40,    0.60,    0.20)
        (0.30,    0.70,    0.20)       
        (0.30,    0.70,    0.40)  
}; \addlegendentry{clay loam}

     \addplot3 coordinates {  
     (0.60,    0.40,    0.40)
     (0.40,    0.60,    0.40)     
     (0.40,    0.60,    0.60)
};  \addlegendentry{dusty clay}

     \addplot3 coordinates {  
      (0.40,    0.60,    0.00)
      (0.40,    0.60,    0.20)
      (0.30,    0.70,    0.20)     
      (0.30,    0.70,    0.00)
}; \addlegendentry{silty clay loam}

     \addplot3 coordinates {  
      (0.15,    0.85,    0.00)
     (0.15,    0.85,    0.15)
     (0.30,    0.70,    0.20)  
     (0.30,    0.70,    0.00)   
};  \addlegendentry{clay dusty loam}

     \addplot3 coordinates {  
     (0.00,    1.00,    0.0)
     (0.00,    1.00,    0.25)
     (0.15,    0.85,    0.15)  
     (0.15,    0.85,    0.00)    
     (0.00,    1.00,    0.0)
};  \addlegendentry{thin loam}

\addplot3 coordinates {  
     (0.30,    0.70,    0.20)
     (0.10,    0.90,    0.40)
      (0.10,    0.90,    0.50)
     (0.20,    0.80,    0.60)
      (0.30,    0.70,    0.40)
}; \addlegendentry{loam}

 \addplot3 coordinates {
       (0.40,    0.60,    0.60)
       (0.40,    0.60,    0.40)    
        (0.30,    0.70,    0.40)           
        (0.20,    0.80,    0.60)       
        (0.20,    0.80,    0.80)
}; \addlegendentry{sandy clay loam}

     \addplot3 coordinates {  
     (0.10,    0.90,    0.40)
     (0.30,    0.70,    0.20)
};  \addlegendentry{dusty loam}

     \addplot3 coordinates {  
     (0.0,    0.0,    0.25)
     (0.0,    0.0,    0.5)
     (0.30,    0.70,    0.20)
     (0.15,    0.85,    0.15)
};    
\node at (axis cs:0.60,0.40,0.2) {\footnotesize{clay}};
\end{ternaryaxis} % <-- finish ternary diagram
\end{tikzpicture} % <-- finish 1st subplot


	\vskip 10pt %  <-- space between two subplots
\begin{tikzpicture} [scale=1]
\begin{ternaryaxis}
[area style,title={\textbf{Granulometric composition of the soil}. Fractions: dust, clay, sand}, %  <-- area style indicate visualization by colors
 legend cell align=left,
 legend pos=outer north east,
 legend entries={sand, loamy sand, sandy loam, sandy clay loam, sandy clay, clay, clay loam, silty clay, silty clay loam, clay loam, silty clay loam, fine loam, silty clay}, 
    xlabel=\footnotesize{Clay (<0.002 mm)}, 
    ylabel=\footnotesize{Dust (0.002-0.05 mm)}, 
    zlabel=\footnotesize{Sand (0.05-2 mm)},
    label style={sloped},
    minor tick num=2] 
   \addplot3 [fill=beaublue] coordinates {  %  <-- we explicitly indicate colors to paint (default: random colors)
         (0.60,    0.40,    0.00)
        (0.40,    0.60,    0.00)
        (0.40,    0.60,    0.20)
};  \addlegendentry{sandy clay})   
      \addplot3 [fill=brilliantrose] coordinates { 
        (0.20,    0.80,    0.80)
        (0.20,    0.80,    0.60)
        (0.10,    0.90,    0.50)
        (0.10,    0.90,    0.40)
       (0.00,    1.00,    0.50)  
        (0.00,    1.00,    0.60)       
        (0.10,    0.90,    0.90)
 };  \addlegendentry{sandy loam}
   \addplot3 [fill=brightturquoise] coordinates { 
         (0.0,    1.00,    0.60) 
        (0.0,    1.00,    0.80)
        (0.10,    0.95,    0.95)
         (0.15,    0.85,    0.85)
};  \addlegendentry{loamy sand}
   \addplot3 [fill=bananayellow] coordinates { 
        (1.00,    0.00,    0.00) 
        (0.60,    0.40,    0.00)        
        (0.40,    0.60,    0.20)    
        (0.40,    0.60,    0.40)
        (0.60,    0.40,    0.40)     
}; \addlegendentry{clay} 
   \addplot3 coordinates { 
      (0.0,    1.00,    1.00) 
      (0.10,    0.85,    0.95) 
      (0.00,    1.00,    0.80)
};  \addlegendentry{sand}
  
     \addplot3 [fill=blue-violet] coordinates {  
        (0.40,    0.60,    0.40)
        (0.40,    0.60,    0.20)
        (0.30,    0.70,    0.20)       
        (0.30,    0.70,    0.40)  
}; \addlegendentry{clay loam}

     \addplot3  [fill=americanrose] coordinates {  
     (0.60,    0.40,    0.40)
     (0.40,    0.60,    0.40)     
     (0.40,    0.60,    0.60)
};  \addlegendentry{dusty clay}

     \addplot3 [fill=applegreen] coordinates {  
      (0.40,    0.60,    0.00)
      (0.40,    0.60,    0.20)
      (0.30,    0.70,    0.20)     
      (0.30,    0.70,    0.00)
}; \addlegendentry{silty clay loam}

     \addplot3 coordinates {  
      (0.15,    0.85,    0.00)
     (0.15,    0.85,    0.15)
     (0.30,    0.70,    0.20)  
     (0.30,    0.70,    0.00)   
};  \addlegendentry{clay dusty loam}

     \addplot3 [fill=amber] coordinates {  
     (0.00,    1.00,    0.0)
     (0.00,    1.00,    0.25)
     (0.15,    0.85,    0.15)  
     (0.15,    0.85,    0.00)    
     (0.00,    1.00,    0.0)
};  \addlegendentry{thin loam}

\addplot3 [fill=Cyan] coordinates {  
     (0.30,    0.70,    0.20)
     (0.10,    0.90,    0.40)
      (0.10,    0.90,    0.50)
     (0.20,    0.80,    0.60)
      (0.30,    0.70,    0.40)
}; \addlegendentry{loam}

 \addplot3 [fill=brilliantlavender] coordinates {
       (0.40,    0.60,    0.60)
       (0.40,    0.60,    0.40)    
        (0.30,    0.70,    0.40)           
        (0.20,    0.80,    0.60)       
        (0.20,    0.80,    0.80)
}; \addlegendentry{sandy clay loam}

     \addplot3 coordinates {  
      (0.30,    0.70,    0.20)
     (0.15,    0.85,    0.15)   
     (0.00,    1.00,    0.25)
     (0.00,    1.00,    0.50)
};  \addlegendentry{dusty loam}

    
    \node at (axis cs:0.60,0.40,0.2) {\footnotesize{clay}};
    
\end{ternaryaxis} % <-- finish ternary
\end{tikzpicture} % <-- finish sublplot 2
	\begin{flushleft} % <--Annotate plot. Align: left.
		\textbf{Granulometric composition of the soil}
	\end{flushleft}% <-- finish plot annotation
\end{document}
